% !TeX encoding = UTF-8
\documentclass[variant=docente,cover=institutional,mode=screen]{ua-engineering-v6-article}
% !TeX encoding = UTF-8
\UASetUniversity{Universidad Austral}
\UASetFaculty{Facultad de Ingeniería}
\UASetDepartment{Departamento de Computación}
\UASetCampus{Pilar}
\UASetCareer{Ingeniería Informática}
\UASetCommission{Comisión A}
\UASetProfessor{Equipo docente}
\UASetSemester{Segundo cuatrimestre 2026}
\UASetCourse{Sistemas Distribuidos}
\UASetDate{2026}


\UASetClassNumber{13}
\UASetClassTitle{Confiabilidad operativa}
\UASetDocKind{Guía docente}

\title{Clase 13 --- Guía docente}
\author{\UAProfessor}
\date{\UADate}

\begin{document}
\maketitle

\section{Objetivo pedagógico profundo}

La meta de esta clase es que el estudiante comprenda que la confiabilidad operativa no se maneja con intuiciones del tipo “esto anda bien” o “esto anda mal”, sino con objetivos formales, métricas significativas y presupuestos de error que permitan gobernar el sistema.

La idea que debe quedar instalada es:

\begin{center}
\textbf{confiabilidad operativa = calidad de servicio administrada explícitamente}
\end{center}

\section{Resultados de aprendizaje esperados}

Al finalizar, el estudiante debería poder:

\begin{itemize}
\item distinguir claramente entre SLI, SLO y SLA;
\item diseñar indicadores orientados al usuario;
\item justificar por qué un objetivo debe ser explícito y medible;
\item calcular e interpretar un error budget;
\item explicar cómo estos conceptos influyen en operación y priorización.
\end{itemize}

\section{Estructura sugerida para 90 minutos}

\subsection*{Bloque 1 (15 min): problema de operar sin objetivos}
Abrir con una pregunta:
\begin{quote}
¿Qué significa realmente que el sistema “ande bien”?
\end{quote}

Buscar respuestas vagas y mostrar su insuficiencia.

\paragraph{Objetivo didáctico}
Instalar la necesidad de formalizar expectativas de servicio.

\subsection*{Bloque 2 (20 min): SLI, SLO y SLA}
Definir cada concepto y trabajar ejemplos buenos y malos.

\paragraph{Objetivo didáctico}
Que el estudiante no vea siglas aisladas, sino una estructura coherente:
métrica → objetivo → acuerdo.

\subsection*{Bloque 3 (20 min): error budgets}
Explicar cómo el error budget transforma el objetivo de confiabilidad en una herramienta de decisión.

\paragraph{Objetivo didáctico}
Mostrar el vínculo entre confiabilidad y velocidad de cambio.

\subsection*{Bloque 4 (15 min): diseño de buenos objetivos}
Discutir:
\begin{itemize}
\item usuario vs infraestructura;
\item promedios vs percentiles;
\item ventanas temporales;
\item objetivos realistas vs imposibles.
\end{itemize}

\paragraph{Objetivo didáctico}
Que el estudiante vea que definir un buen SLO es un problema de diseño, no de burocracia.

\subsection*{Bloque 5 (20 min): operación guiada por objetivos}
Mostrar cómo cambia la toma de decisiones:
\begin{itemize}
\item incidentes;
\item despliegues;
\item deuda técnica;
\item conflicto producto/ingeniería.
\end{itemize}

\paragraph{Objetivo didáctico}
Conectar métricas con gobierno real del sistema.

\section{Qué explicar y por qué}

\subsection*{1. SLI}
Porque es la base medible de todo el esquema.

\subsection*{2. SLO}
Porque convierte expectativas vagas en objetivos explícitos.

\subsection*{3. SLA}
Porque distingue claramente objetivos internos de compromisos externos.

\subsection*{4. Error budget}
Porque introduce la lógica económica y operativa de la confiabilidad.

\subsection*{5. Operación guiada por objetivos}
Porque muestra que estas ideas no son meramente descriptivas: transforman cómo se decide.

\section{Preguntas disparadoras recomendadas}

\begin{itemize}
\item ¿Qué parte del servicio importa realmente al usuario?
\item ¿Qué métrica se parece a salud del sistema, pero no a salud del servicio?
\item ¿Qué significa que un objetivo sea demasiado estricto o demasiado laxo?
\item ¿Qué harías si el error budget se consume demasiado rápido?
\item ¿Qué conflicto resuelve un SLO entre producto e ingeniería?
\end{itemize}

\section{Errores frecuentes del alumnado}

\begin{itemize}
\item confundir SLI con cualquier métrica disponible;
\item definir objetivos internos imposibles de sostener;
\item usar promedios cuando el problema real vive en percentiles;
\item tratar el error budget como permiso para fallar sin criterio;
\item no distinguir objetivo técnico de acuerdo contractual.
\end{itemize}

\section{Ejercicio en vivo sugerido}

Tomar tres servicios:
\begin{itemize}
\item checkout;
\item API pública;
\item pipeline de eventos.
\end{itemize}

Pedir al curso que proponga para cada uno:
\begin{itemize}
\item un SLI relevante;
\item un SLO razonable;
\item qué error budget implicaría;
\item qué decisión tomarían si ese budget se consume demasiado rápido.
\end{itemize}

\paragraph{Objetivo}
Que el estudiante pase de repetir definiciones a diseñar objetivos concretos y usables.

\section{Momento crítico de la clase}

El momento importante ocurre cuando el estudiante entiende que:

\begin{center}
\textbf{la confiabilidad operativa no exige 100\% de perfección, sino un nivel de servicio explícito y gobernable}
\end{center}

Ese punto debe remarcarse, porque corrige una intuición muy frecuente: creer que la única confiabilidad valiosa es la perfección absoluta.

\section{Criterio de éxito docente}

La clase salió bien si el estudiante puede:

\begin{itemize}
\item proponer un SLI relevante;
\item convertirlo en un SLO razonable;
\item explicar qué error budget resulta;
\item usar ese marco para justificar decisiones de operación o despliegue.
\end{itemize}

\section{Conexión con la clase siguiente}

Esta clase deja abierta una transición natural hacia:
\begin{itemize}
\item gestión de incidentes;
\item postmortems;
\item aprendizaje organizacional;
\item cultura operativa.
\end{itemize}

\begin{uakeybox}
Pregunta puente: si ya sabemos qué significa “cumplir servicio”, ¿cómo respondemos, aprendemos y mejoramos cuando inevitablemente el sistema no cumple?
\end{uakeybox}


\section{Plan de conducción ampliado}
La clase debe alternar explicación, predicción y contraste. Antes de revelar cada mecanismo, pedir al grupo que anticipe qué puede salir mal; después, volver al mismo escenario y preguntar qué propiedad mejora y qué costo aparece. Cada bloque conceptual debe cerrar con una decisión de diseño observable.
\subsection*{Secuencia recomendada}
\begin{enumerate}
\item Recuperar el prerequisito de la clase anterior.
\item Presentar un caso mínimo que falle y solicitar hipótesis.
\item Formalizar el concepto central de Confiabilidad operativa, distinguiendo seguridad, progreso y supuestos.
\item Resolver un ejemplo completo en pizarra, incluyendo una falla intermedia.
\item Comparar una alternativa de diseño y explicitar trade-offs.
\item Cerrar con una pregunta del Trabajo Final y una pregunta de defensa oral.
\end{enumerate}
\section{Diagnóstico de comprensión durante la clase}
No preguntar solamente ``¿se entendió?''. Usar preguntas discriminantes: ``¿qué cambia si el mensaje llega dos veces?'', ``¿qué propiedad sigue siendo cierta si cae un nodo?'', ``¿esto demuestra seguridad o progreso?'', ``¿qué observa un cliente externo?''. Si el estudiante responde solo con nombres de patrones, pedir una ejecución concreta.
\section{Criterios para corregir ejercicios}
Una respuesta excelente declara supuestos, construye una secuencia verificable, nombra la garantía con precisión, reconoce un costo y conecta la decisión con el dominio. Penalizar afirmaciones absolutas sin condiciones.
\section{Vinculación con el Trabajo Final Integrador}
Reservar entre 5 y 10 minutos para que cada grupo registre una decisión con: problema, alternativa elegida, alternativa descartada y razón. Esto crea trazabilidad entre las clases y las entregas acumulativas.
\section{Lectura docente previa}
Revisar la bibliografía indicada en los apuntes y preparar al menos un contraejemplo que no aparezca literalmente en las diapositivas. El objetivo es poder responder preguntas de frontera sin convertir la clase en una lectura del deck.

\section{Arquitectura didáctica v9 para 180 minutos}
\subsection*{Bloque 1 - desafío inicial (20 min)}
Presentar una situación donde aparezca SLI mal elegido, ventana incorrecta, budget consumido y alerta ruidosa. Pedir predicciones individuales antes de introducir vocabulario. El objetivo es hacer visibles los supuestos intuitivos y recuperar el prerequisito de la clase anterior.

\subsection*{Bloque 2 - construcción conceptual (35 min)}
Formalizar SLI, SLO, SLA, error budgets, burn rates y operación guiada por objetivos. Separar seguridad, progreso y experiencia del usuario. Explicar no solo \emph{qué} hace cada mecanismo, sino \emph{por qué} existe y qué supuesto permite que funcione.

\subsection*{Bloque 3 - modelo y figura (35 min)}
Construir en pizarra una línea temporal, grafo, log, máquina de estados o tabla de decisiones. Recorrer una ejecución nominal y una adversarial. La representación debe quedar como referencia para los apuntes y el ejercicio principal.

\subsection*{Pausa (10 min)}
Recoger una pregunta anónima o un punto de confusión. Elegir una pregunta que permita distinguir síntoma, causa y propiedad.

\subsection*{Bloque 4 - caso trabajado con falla (40 min)}
Aplicar el mecanismo al caso conductor. Introducir una falla a mitad de ejecución y detenerse en cada transición: quién sabe qué, qué se persiste, qué garantía sigue vigente y qué observa el cliente.

\subsection*{Bloque 5 - taller de decisión (25 min)}
Los grupos aplican P-F-G-S-T a su Trabajo Final. Deben producir un ADR breve con alternativa elegida, alternativa descartada y evidencia pendiente. La evidencia de esta clase es SLI calculado, SLO, budget, burn rate y política de decisión.

\subsection*{Bloque 6 - cierre y exit ticket (15 min)}
Cada estudiante responde: propiedad principal, supuesto decisivo, costo aceptado y condición bajo la cual cambiaría de diseño. Usar las respuestas para iniciar la clase siguiente.

\section{Qué explicar, por qué y cómo verificarlo}
\begin{itemize}
\item Explicar el problema antes del producto, porque reconocer la familia de problema es el resultado cognitivo central.
\item Explicar el invariante, porque permite evaluar configuraciones y no solo repetir un flujo nominal.
\item Explicar el límite, porque SLI mal elegido, ventana incorrecta, budget consumido y alerta ruidosa puede invalidar supuestos o reducir progreso.
\item Verificar comprensión mediante SLI calculado, SLO, budget, burn rate y política de decisión, no mediante una definición aislada.
\end{itemize}

\section{Preguntas socráticas}
\begin{itemize}
\item ¿Qué sabe cada actor en este punto de la ejecución?
\item ¿Qué propiedad sigue siendo cierta si la respuesta nunca llega?
\item ¿Qué parte es safety y cuál es liveness?
\item ¿Qué alternativa más simple resolvería el problema si relajamos una garantía?
\item ¿Qué observación refutaría nuestra hipótesis de diseño?
\end{itemize}

\section{Errores previsibles y estrategia de corrección}
\begin{itemize}
\item Si el alumno responde con un nombre de tecnología, pedir actores, mensajes y estados.
\item Si confunde timeout con falla, construir dos ejecuciones indistinguibles.
\item Si promete todas las garantías, pedir comportamiento bajo partición o saturación.
\item Si mide solo rendimiento, preguntar cómo evidencia la propiedad semántica.
\item Si la solución es excesiva, pedir la alternativa mínima y el costo marginal de la garantía fuerte.
\end{itemize}

\section{Criterios de evaluación formativa}
Una respuesta excelente declara supuestos, recorre una ejecución, nombra con precisión la garantía, reconoce el costo y propone evidencia reproducible. Una respuesta suficiente identifica el mecanismo y el trade-off principal. Una respuesta insuficiente usa slogans, omite fallas o confunde producto con propiedad.

\section{Preparación docente y bibliografía}
Lectura previa recomendada: Google SRE y SRE Workbook, capítulos de SLOs y error budgets. Preparar un contraejemplo adicional y una variante del caso en la que cambie el costo del error; esto evita convertir la clase en una lectura de diapositivas.

\end{document}
